%%-----------Chapter 5------------------------------------------
\chapter{Further Applications of Newton's Laws}

% ==============================================================
\section{Contact Forces}
% ==============================================================

Pushing, lifting and pulling are \textbf{\textit{contact forces}} that we experience in the everyday world. Rest your hand on a table; the atoms that form the molecules that make up the table and your hand are in contact with each other. If you press harder, the atoms are also pressed closer together. The electrons in the atoms begin to repel each other and your hand is pushed in the opposite direction by the table.

According to Newton's Third Law, the force of your hand on the table is equal in magnitude and opposite in direction to the force of the table on your hand. Clearly, if you push harder the force increases. If you push your hand straight down on the table, the table pushes back in a direction perpendicular (normal) to the surface. Slide your hand gently forward along the surface of the table. You barely feel the table pushing upward, but you do feel the friction acting as a resistive force to the motion of your hand. This force acts tangential to the surface and opposite to the motion of your hand. Push downward and forward. Try to estimate the magnitude of the force acting on your hand.

The force of the table acting on your hand, $\vec{F}_{c} \equiv \vec{C}$, is called the \textbf{\textit{contact force}}. This force has both a normal component to the surface, $\vec{C}_{\perp} \equiv \vec{\mathscr{N}}$, called the \textbf{\textit{normal force}}, and a tangential component to the surface, $\vec{C}_{\parallel} \equiv \vec{f}$, called the \textbf{\textit{friction force}} (Figure~\ref{fig:contactforce}).



\begin{figure}[ht]
\centering
\begin{tikzpicture}[>=Stealth, thick, scale=1.2]
    % --- Warm paper background ---
    \fill[paper, rounded corners=6pt] (-3.0,-0.50) rectangle (5.4,3.0);
    
    % Surface parallelogram
    \fill[gray!12] (-2.5,-0.3) -- (3.5,-0.3) -- (4.5,0.7) -- (-1.5,0.7) -- cycle;
    \draw[gray!50, thin] (-2.5,-0.3) -- (3.5,-0.3) -- (4.5,0.7) -- (-1.5,0.7) -- cycle;
    % Block
    \fill[gray!25] (0.65, 0.2) rectangle (1.75, 1.15);
    \draw[thick] (0.65, 0.2) rectangle (1.75, 1.15);
    % Contact origin
    \coordinate (O) at (1.2, 0.2);
    % F: red, pushing left
    \draw[->, very thick, myred] (0.65, 0.67) -- ++(-1.5, 0)
        node[left, font=\small] {$\vec{F}$};
    % N: blue, up
    \draw[->, thick, myblue] (O) -- ++(0, 2.2)
        node[right, font=\small] {$\vec{C}_\perp \equiv \vec{\mathscr{N}}$};
    % f: green, horizontal right
    \draw[->, thick, mygreen] (O) -- ++(2.2, 0)
        node[right, font=\small] {$\vec{C}_{\parallel} \equiv \vec{f}$};
    % C: black, resultant
    \draw[->, very thick, black] (O) -- ++(2.2, 2.2)
        node[above right, font=\small] {$\vec{C}$};
    % Dashed construction lines
    \draw[dashed, gray!55, thin] (O) ++(0, 2.2) -- ++(2.2, 0);
    \draw[dashed, gray!55, thin] (O) ++(2.2, 0) -- ++(0, 2.2);
\end{tikzpicture}
\caption{The block is pushed to the left by $\vec{F}$. The contact force $\vec{C}$ decomposes into a normal component $\vec{\mathscr{N}}$ and a tangential friction component $\vec{f}$ opposing the motion.}
\label{fig:contactforce}
\end{figure}



The contact force, written in terms of its component forces, is therefore
\begin{equation}
    \vec{C} = \vec{C}_\perp + \vec{C}_{\parallel} \equiv \vec{\mathscr{N}} + \vec{f}\,.
\end{equation}
The normal component $\vec{\mathscr{N}}$ and the tangential component $\vec{f}$ are not independent forces but the vector components of the contact force, perpendicular and parallel to the surface of contact. 



% ==============================================================
\section*{Friction}
% ==============================================================

When a block slides along a horizontal surface or down an inclined plane, a lateral force resists the motion. Even when the block is at rest on an inclined plane, a lateral force prevents it from sliding. This resistive force is known as \textit{dry friction}, and
there are two distinguishing types when surfaces are in contact with each other.

When two surfaces move relative to one another, the friction is called \textbf{\textit{kinetic friction}} (or \textbf{\textit{sliding friction}}). When the surfaces are stationary relative to each other but a lateral force is present, the friction is called \textbf{\textit{static friction}}.

\begin{predictbox}[5.1: The brick, flat or standing?]
Take a brick (or this textbook). You can drag it lying \emph{flat} (large face
down) or standing on its \emph{edge} (small face down). The weight is the same
either way; only the contact area changes --- by a factor of three or four.
Predict: to slide it at steady speed, which orientation needs the harder
pull? Flat, on edge, or no difference? Write your answer down with a one-line
reason --- then, if a spring scale or a rubber band is handy, test it before
reading on.
\end{predictbox}

You have just repeated, with a brick, the experiments Leonardo da Vinci
recorded in his notebooks over the years 1493 to about 1515 --- the first
systematic measurements of kinetic friction ever made. Dragging blocks with
hanging weights and pulleys, he found the two facts that still carry his
fingerprints today: the force needed to slide a block is proportional to how
hard the surfaces are pressed together, and --- the surprise --- it does
\emph{not} depend on the area of contact. Da Vinci never published; his
results were rediscovered and published by Guillaume Amontons in 1699, and
extended by Charles-Augustin Coulomb, who showed that kinetic friction is
also independent of sliding speed for ordinary speeds. Hold on to your brick
prediction; the second of da Vinci's facts will grade it in a moment.

\subsection{Kinetic Friction}

The magnitude of the kinetic friction force $\vec{f}^{\,k}$ is proportional to the normal force between the two surfaces,
\begin{equation}
    f_k = \mu_k \mathscr{N}\,,
\end{equation}
where $\mu_k$ is the \textbf{\textit{coefficient of kinetic friction}}, a dimensionless constant depending only on the two materials in contact. Now look at what this law does \emph{not} contain: the contact area. The flat brick and the standing brick need exactly the same pull --- if your prediction in Stop and Predict 5.1 said otherwise, you are with the majority, and with common sense; but the experiment (da Vinci's and yours) says no difference. Two blocks of the same mass but different footprints require the same force to slide at constant speed (Figure~\ref{fig:kineticarea}). Roughly speaking: the flat brick spreads its weight over more contact points, but presses each one more lightly; the standing brick concentrates the same weight on fewer, harder-pressed points. The two effects cancel. This follows from the microscopic structure of surfaces, as discussed in the next section.

\begin{figure}[ht]
\centering
\begin{tikzpicture}[>=Stealth, thick]
    % --- Warm paper background ---
    \fill[paper, rounded corners=6pt] (-6.5,-0.9) rectangle (6,1.7);
    \begin{scope}[shift={(-3.6,0)}]
        \draw[thick] (-1.55,-0.02) -- (1.55,-0.02);
        \fill[pattern=north east lines, pattern color=gray!45]
            (-1.55,-0.22) rectangle (1.55,-0.02);
        \draw[fill=gray!15] (-1.2,0) rectangle (1.2,0.72);
        \draw[->, thick] (-1.2,0.36) -- ++(-1.1,0)
            node[left, font=\small] {$\vec{F}$};
        \draw[->, thick] ( 1.2,0.36) -- ++( 1.1,0)
            node[right, font=\small] {$\vec{f}_k$};
        \node[font=\small\itshape] at (0,-0.55) {(a)};
    \end{scope}
    \begin{scope}[shift={(3.6,0)}]
        \draw[thick] (-1.0,-0.02) -- (1.0,-0.02);
        \fill[pattern=north east lines, pattern color=gray!45]
            (-1.0,-0.22) rectangle (1.0,-0.02);
        \draw[fill=gray!15] (-0.5,0) rectangle (0.5,1.44);
        \draw[->, thick] (-0.5,0.72) -- ++(-1.1,0)
            node[left, font=\small] {$\vec{F}$};
        \draw[->, thick] ( 0.5,0.72) -- ++( 1.1,0)
            node[right, font=\small] {$\vec{f}_k$};
        \node[font=\small\itshape] at (0,-0.55) {(b)};
    \end{scope}
\end{tikzpicture}
\caption{Kinetic friction is independent of contact area. Block~(a) has twice the footprint of block~(b), yet both require the same force to slide at constant speed.}
\label{fig:kineticarea}
\end{figure}

\subsection{Static Friction}

Static friction does not take a fixed value; it adjusts itself to prevent relative motion. As you push a stationary block with increasing force, static friction matches your push exactly until a maximum value is reached. Beyond this threshold, the block slips and kinetic friction takes over. The magnitude of static friction therefore satisfies
\begin{equation}
    0 \leq f_s \leq (f_s)_{\max}\,,
\end{equation}
where the maximum is found experimentally to be proportional to the normal force,
\begin{equation}
    (f_s)_{\max} = \mu_s \mathscr{N}\,.
\end{equation}
Here $\mu_s$ is the \textbf{\textit{coefficient of static friction}}. Experiment consistently gives $\mu_s > \mu_k$: more force is needed to initiate sliding than to sustain it. This small difference accounts for the stick-slip behaviour of chalk on a blackboard, fingernails on glass, or a violin bow on a string.

The direction of static friction always opposes the applied force, as shown in Figure~\ref{fig:staticdir}.

\begin{figure}[ht]
\centering
\begin{tikzpicture}[>=Stealth, thick]
    \fill[paper, rounded corners=6pt] (-6.0,-0.9) rectangle (6,1.2);
    \begin{scope}[shift={(-3.3,0)}]
        \draw[thick] (-1.4,-0.02) -- (1.4,-0.02);
        \fill[pattern=north east lines, pattern color=gray!45]
            (-1.4,-0.22) rectangle (1.4,-0.02);
        \draw[fill=gray!15] (-0.85,0) rectangle (0.85,0.9);
        \draw[->, thick] (-0.85,0.45) -- ++(-1.1,0)
            node[left,  font=\small] {$\vec{F}$};
        \draw[->, thick] ( 0.85,0.45) -- ++( 1.0,0)
            node[right, font=\small] {$\vec{f}_s$};
        \node[font=\small\itshape] at (0,-0.55) {(a)};
    \end{scope}
    \begin{scope}[shift={(3.3,0)}]
        \draw[thick] (-1.4,-0.02) -- (1.4,-0.02);
        \fill[pattern=north east lines, pattern color=gray!45]
            (-1.4,-0.22) rectangle (1.4,-0.02);
        \draw[fill=gray!15] (-0.85,0) rectangle (0.85,0.9);
        \draw[->, thick] ( 0.85,0.45) -- ++( 1.1,0)
            node[right, font=\small] {$\vec{F}$};
        \draw[->, thick] (-0.85,0.45) -- ++(-1.0,0)
            node[left,  font=\small] {$\vec{f}_s$};
        \node[font=\small\itshape] at (0,-0.55) {(b)};
    \end{scope}
\end{tikzpicture}
\caption{Static friction opposes the applied force. In~(a) $\vec{F}$ acts to the left so $\vec{f}_s$ acts to the right. In~(b) the roles are reversed.}
\label{fig:staticdir}
\end{figure}

Two important distinctions between static and kinetic friction are worth keeping in mind.
\vspace{-0.1cm}
\begin{enumerate}[leftmargin=2em, itemsep=2pt]
    \item The magnitude and direction of static friction depend on the applied forces and adjust to maintain equilibrium, up to the maximum $(f_s)_{\max}$. Kinetic friction, by contrast, has a fixed magnitude $\mu_k \mathscr{N}$ for a given normal force.
    \item When the applied force along the contact surface reaches $(f_s)_{\max}$, the surfaces begin to slip. This threshold condition is called the \textbf{\textit{just slipping}} condition, and marks the transition from static to kinetic friction.
\end{enumerate}

\vspace{0.3cm}

\begin{workedexample}[5.1: Sliding Block]
A block of mass \(m\) slides down an incline of angle \(30^\circ\) with an acceleration \(g/4\). Find the kinetic friction coeffcient \(\mu_k\).


\medskip
\textit{Solution.}
We resolve the forces into components parallel and perpendicular to the incline, taking
the positive direction down the slope for motion and outward from the surface for the
normal direction. 

\begin{center}
\begin{tikzpicture}[scale=1.7, line width=0.9pt, >=stealth]

    \pgfmathsetmacro{\ang}{30}

    % --- Warm paper background ---
    \fill[paper, rounded corners=6pt] (-0.7,-0.55) rectangle (5.9,3.5);

    % --- Incline triangle ---
    \coordinate (A) at (0,0);
    \coordinate (B) at (5,0);
    \coordinate (C) at (0, {5*tan(\ang)});
    \draw[thick, fill=gray!12] (A) -- (B) -- (C) -- cycle;

    % Ground hatch
    \draw[thick] (-0.2,0) -- (5.25,0);
    \foreach \x in {-0.1,0.3,0.7,1.1,1.5,1.9,2.3,2.7,3.1,3.5,3.9,4.3,4.7,5.1}
        \draw[thin] (\x,0) -- +(-0.18,-0.14);

    % Angle at B
    \draw (4.2,0) arc (180:{180-\ang}:0.8);
    \node at (3.9,0.28) {\small $30^\circ$};

    % --- Block position on slope ---
    \pgfmathsetmacro{\bx}{5 - 0.50*5}
    \pgfmathsetmacro{\by}{0.50*5*tan(\ang)}
    \pgfmathsetmacro{\bcx}{\bx + 0.275*cos(60)}
    \pgfmathsetmacro{\bcy}{\by + 0.275*sin(60)}
    \coordinate (O) at (\bcx,\bcy);

    \begin{scope}[shift={(O)}, rotate={-\ang}]
        \draw[fill=white, thick] (-0.35,-0.275) rectangle (0.35,0.275);
    \end{scope}

    % --- Lengths (shortened) ---
    \pgfmathsetmacro{\L}{1.05}
    \pgfmathsetmacro{\Lsin}{\L*sin(\ang)}
    \pgfmathsetmacro{\Lcos}{\L*cos(\ang)}

    \coordinate (mgTip)   at ($(O)+(270:\L)$);
    \coordinate (sinTip)  at ($(O)+(-\ang:\Lsin)$);
    \coordinate (cosTip)  at ($(O)+(240:\Lcos)$);

    % --- Dashed parallelogram of components ---
    \draw[dashed, gray!55] (sinTip) -- (mgTip);
    \draw[dashed, gray!55] (cosTip) -- (mgTip);

    % --- Weight and components ---
    \draw[->, thick, black] (O) -- (mgTip)
        node[below, font=\small] {$mg$};
    \draw[->, thick, black!55] (O) -- (sinTip)
        node[below right, font=\scriptsize, xshift=2pt, yshift=-1pt] {$mg\sin\theta$};
    \draw[->, thick, black!55] (O) -- (cosTip)
        node[below left, font=\scriptsize] {$mg\cos\theta$};

    % --- angle theta between mg and mg cos th ---
    \draw[line width=0.5pt, gray!70] (O)+(270:0.40) arc (270:240:0.40);
    \node[font=\scriptsize, gray!40!black] at ($(O)+(255:0.55)$) {$\theta$};

    % --- Normal force ---
    \draw[->, thick, blue!70!black] (O) -- ($(O)+(60:\Lcos)$)
        node[above right, font=\small, xshift=-2pt] {$\mathscr{N}$};

    % --- Friction ---
    \draw[->, thick, orange!80!black] (O) -- ($(O)+(150:\Lsin)$)
        node[above left, font=\small, xshift=2pt] {$f_k$};

    % --- Acceleration g/4 ---
    \coordinate (aStart) at ($(O)+(60:0.43)+(-30:0.30)$);
    \draw[->, very thick, green!39!black] (aStart) -- ++(-\ang:0.95)
        node[above right, font=\small, xshift=-3pt] {$a=g/4$};

    % --- Coordinate axes inset ---
    \coordinate (Oax) at (4.55,2.45);
    \draw[->, gray!55!black, line width=0.8pt]
        (Oax) -- ($(Oax)+(-\ang:0.7)$) node[right, font=\scriptsize, gray!45!black] {$x$};
    \draw[->, gray!55!black, line width=0.8pt]
        (Oax) -- ($(Oax)+(60:0.7)$) node[above, font=\scriptsize, gray!45!black] {$y$};
    \draw[gray!55!black, line width=0.5pt]
        ($(Oax)+(-\ang:0.16)$) -- ($(Oax)+(-\ang:0.16)+(60:0.16)$) -- ($(Oax)+(60:0.16)$);

\end{tikzpicture}
\end{center}


\vspace{0.3cm}
Perpendicular to the incline there is no acceleration, so the normal force balances the
perpendicular component of gravity,
\begin{equation*}
    \mathscr{N} - mg\cos\theta = 0 \implies \mathscr{N} = mg\cos\theta.
\end{equation*}

Along the incline, the block accelerates down the slope at $a = g/4$. The driving force is
the parallel component of gravity, $mg\sin\theta$, opposed by kinetic friction
$f_k = \mu_k\mathscr{N}$ acting up the slope. Newton's second law gives
\begin{equation*}
    mg\sin\theta - \mu_k\mathscr{N} = ma.
\end{equation*}
Substituting $\mathscr{N} = mg\cos\theta$ and $a = g/4$,
\begin{equation*}
    mg\sin\theta - \mu_k mg\cos\theta = \tfrac{1}{4}mg.
\end{equation*}
The mass $m$ and the acceleration $g$ cancel throughout, leaving a relation among pure
numbers,
\begin{equation*}
    \sin\theta - \mu_k\cos\theta = \tfrac{1}{4}.
\end{equation*}
Solving for the coefficient of kinetic friction,
\begin{equation*}
    \mu_k = \frac{\sin\theta - \tfrac{1}{4}}{\cos\theta}.
\end{equation*}
With $\theta = 30^\circ$, so that $\sin\theta = \tfrac{1}{2}$ and $\cos\theta = \tfrac{\sqrt{3}}{2}$,
\begin{equation*}
    \mu_k = \frac{1/2 - 1/4}{\sqrt{3}/2}
          = \frac{1}{2\sqrt{3}}.
\end{equation*}



\end{workedexample}


\newpage


\section*{The Normal Force}

When two solid bodies are pressed together, their surfaces resist interpenetration.
Push your hand against a wall and the wall pushes back, preventing your hand from
passing through (Figure~\ref{fig:stickwall}). This perpendicular push is called the
\textit{normal force}, where \textit{normal} means perpendicular to the surface of
contact. Notice that the harder you push on the wall with a force $\vec{F}$, the harder
the wall pushes back on your hand with the normal force $\vec{\mathscr{N}}$.

The origin of this force is atomic. The atoms at the surface of your hand and the
atoms at the surface of the wall repel one another the moment they come into contact,
opposing any further compression. The material of the wall behaves, at the microscopic
level, like an array of tiny springs.\footnote{The ``atomic springs'' picture is a
useful analogy but not the complete story. The true origin of this repulsion is quantum
mechanical: as electron clouds overlap, the Pauli exclusion principle forces electrons
into higher energy states, giving rise to a repulsive force. Classical electrostatics
plays a secondary role. We will encounter the Pauli exclusion principle when we study
modern physics.} When you push against the wall, these atomic springs compress
slightly, and the repulsive force they exert grows with the compression until it exactly
balances your push. For a hard material such as concrete, the compression is so small
as to be invisible, and the wall appears perfectly rigid.

This is why the normal force is not a fixed quantity: it adjusts itself to match whatever
force is pressing the surfaces together. Push harder, and the atomic springs compress
more, increasing the normal force in response. As the figure shows, $\vec{F}$ and
$\vec{\mathscr{N}}$ always grow together and remain equal in magnitude. The normal force is
therefore a \textit{constraint force}, one whose magnitude is determined not by a fixed
law but by the condition that the two surfaces do not pass through each other.


\begin{figure}[ht!]
\centering
\begin{tikzpicture}[scale=1.0, line cap=round, line join=round, >=Stealth]

    \tikzset{limb/.style={line width=2.2pt, black}}

    % --- Warm paper background ---
    \fill[paper, rounded corners=6pt] (-1.55,-0.45) rectangle (6.45,4.15);

    % --- Ground ---
    \draw[line width=1pt, gray!60!brown] (-0.3,0) -- (6.2,0);
    \foreach \x in {0.2,0.7,1.2,4.8,5.3,5.8} {
        \draw[gray!55!brown, line width=0.6pt] (\x,-0.04) -- (\x-0.18,-0.18);
    }

    % --- Wall ---
    \fill[pattern=north east lines, pattern color=gray!45!brown]
        (-0.35,0) rectangle (0,3.9);
    \draw[line width=1.6pt] (0,0) -- (0,3.9);

    % ===== Stick figure =====
    \coordinate (hand)     at (0.05, 2.55);
    \coordinate (shoulder) at (2.05, 2.20);
    \coordinate (hip)      at (3.30, 1.05);
    \coordinate (kneeF)    at (3.05, 0.55);
    \coordinate (footF)    at (3.55, 0.05);
    \coordinate (footB)    at (5.10, 0.05);

    \draw[limb] (hand) -- (shoulder);
    \draw[limb] (0.05,2.40) -- (0.05,2.72);
    \draw[limb] (shoulder) -- (hip);
    \draw[limb] (hip) -- (footB);
    \draw[limb] (footB) .. controls (5.35,0.02) and (5.45,0.10) .. (5.40,0.20);
    \draw[limb] (hip) -- (kneeF) -- (footF);
    \draw[limb] (footF) .. controls (3.80,0.02) and (3.90,0.10) .. (3.85,0.22);

    % --- Head with strained expression ---
    \draw[limb, fill=paper] (2.55,2.62) circle (0.46);
    \draw[line width=1.3pt] (2.30,2.84) -- (2.46,2.72);
    \draw[line width=1.3pt] (2.78,2.80) -- (2.62,2.70);
    \draw[line width=1.1pt] (2.34,2.64) -- (2.45,2.60);
    \draw[line width=1.1pt] (2.74,2.61) -- (2.63,2.58);
    \draw[line width=1.1pt, fill=paper] (2.38,2.30) rectangle (2.74,2.44);
    \draw[line width=0.7pt] (2.50,2.30) -- (2.50,2.44);
    \draw[line width=0.7pt] (2.62,2.30) -- (2.62,2.44);

    % --- Force arrows, both anchored at the hand ---
    \draw[->, line width=1.8pt, myred]
        (0.05,2.6) -- ++(-1.05,0) node[above, font=\small] {$\vec{F}$};
    \draw[->, line width=1.8pt, myblue]
        (0.05,2.6) -- ++(1.05,0) node[above, font=\small] {$\vec{\mathscr{N}}$};
    \fill[black] (hand) circle (1.5pt);

\end{tikzpicture}
\caption{When you push against a wall ($\vec{F}$), the wall pushes back on your hand
with an equal and opposite normal force ($\vec{\mathscr{N}}$), perpendicular to its surface.
Both forces act at the point of contact.}
\label{fig:stickwall}
\end{figure}



\section*{The Normal Force and Weight}

We say we are measuring our weight when we step onto a bathroom scale. But what is the scale actually measuring?

Two forces act on a person standing on a scale (Figure~\ref{fig:scale}). The first is the gravitational force $\vec{F}_g = m\vec{g}$, directed downward. The second is the contact force from the scale; because the person is at rest, this contact force is purely normal and we write it as $\vec{\mathscr{N}}$, directed upward. Since the person is at rest, the vertical acceleration is zero, and Newton's Second Law gives
\begin{equation}
    \vec{\mathscr{N}} + \vec{F}_g = \vec{0}.
\end{equation}
Taking the positive direction upward, this becomes
\begin{equation}
    \mathscr{N} - mg = 0 \implies \mathscr{N} = mg.
\end{equation}

It is worth being precise about four quantities that are easy to confuse. The \textit{true mass} $m$ is the amount of matter in the person, an invariant property that does not change with motion or location.\footnote{Newtonian mechanics treats mass as a fixed quantity. Special relativity shows that an object's inertia grows with its speed. One common way to express this
introduces a velocity-dependent mass $m = m_0/\sqrt{1 - v^2/c^2}$, where $m_0$ is the
rest mass and $c$ is the speed of light; for everyday speeds $v \ll c$ this is
indistinguishable from $m_0$. Modern physics usually prefers instead to keep mass as
the invariant $m_0$ and describe the growing inertia through momentum and energy. Either
way, the effect is utterly negligible at ordinary speeds. We will return to this when we
study modern physics.} The \textit{true weight} $mg$ is the actual gravitational force the Earth exerts on the person. The scale, however, does not measure either of these directly. It measures the normal force $\mathscr{N}$ pressing on it, which we call the \textit{apparent weight}. Dividing by $g$, it displays a number in kilograms, the \textit{apparent mass} $\mathscr{N}/g$. At rest the apparent mass equals the true mass and the distinction seems academic. The elevator, discussed in the next section, will reveal the difference.

The word ``weight'' is often used loosely to mean the gravitational force the Earth exerts on an object. We shall always call that force the \textit{gravitational force} rather than ``weight'', to keep it distinct from what a scale reads. When you jump into the air you feel ``weightless'' because no normal force acts on you, even though the Earth still pulls you down. You have not turned gravity off.

A common misconception deserves to be addressed before we proceed. The normal force and the gravitational force are two completely different forces. They are equal in magnitude here only because the person is at rest, and this equality will fail the moment there is acceleration. More importantly, they do not form a Newton's Third Law interaction pair: both act on the same object (the person), whereas Third Law pairs always act on \textit{different} objects. The true Third Law partner of the gravitational force on the person due to the Earth, $\vec{F}^{\,g}_{E,P}$, is the gravitational force on the Earth due to the person,  $\vec{F}^{\,g}_{P,E}$, and they satisfy \(\vec{F}^{\,g}_{E,P} = -\vec{F}^{\,g}_{P,E}\). These act on different objects: one on the person, the other on the Earth. Likewise, the Third Law partner of the normal force $\vec{\mathscr{N}}$ from the scale is the force the person's feet exert downward on the scale surface.


\section*{Experience in an Elevator}

Now we take the bathroom scale somewhere interesting. If a scale measures
the normal force --- not gravity --- then the reading should change whenever the
normal force has to do something other than just hold you up. An elevator is
the perfect laboratory: it can accelerate up, accelerate down, cruise, and
(in a thought experiment we will keep hypothetical) drop freely.

\begin{predictbox}[5.2: Four rides on a scale]
A 60\,kg person stands on a scale in an elevator. At rest the scale reads
60.0\,kg. Predict --- more than 60, less than 60, exactly 60, or zero? --- for
each ride, \emph{before} reading the analysis:
\begin{enumerate}[leftmargin=2em, itemsep=1pt]
  \item The elevator moves \emph{upward at a steady} $5\,\mathrm{m/s}$.
  \item The elevator \emph{accelerates upward}.
  \item The elevator \emph{accelerates downward} (slowing while going up, or
        speeding up while going down).
  \item The cable snaps and the elevator \emph{falls freely}.
\end{enumerate}
Ride 1 is the one people argue about longest --- ``moving up'' feels like it
should read more. Commit to all four; the next pages check them one by one.
\end{predictbox}

\noindent\textit{On the companion site:} the \textbf{elevator simulator}
(\texttt{Chapter 5} on the course website) has an acceleration slider and a
live scale readout --- test your four predictions there before or after the
analysis, and watch the cable snap in case~4.

In each case exactly two forces act on the person: the gravitational force $m\vec{g}$
directed downward and the normal force $\vec{\mathscr{N}}$ directed upward.
 
\vspace{0.3cm}
\noindent\textbf{\sffamily(a)\; At Rest or Moving at Constant Velocity.}
 
\begin{minipage}[t]{0.55\textwidth}
\vspace{0pt}
When the elevator is stationary or moving at constant velocity, the acceleration is zero.
The person's state of motion is not changing, so the floor need only support the person's
weight, nothing more. Newton's second law in the vertical direction gives
\begin{equation*}
    \vec{\mathscr{N}} + \vec{F}_g = \vec{0}
    \implies \mathscr{N} - mg = 0,
\end{equation*}
so that
\begin{equation*}
    \mathscr{N} = mg = (60)(10) = 600\;\mathrm{N}.
\end{equation*}
The scale displays $\mathscr{N}/g = 60.0\,\mathrm{kg}$. Here the apparent weight equals the
true weight, which is why we never notice the distinction in ordinary standing.
Note what this settles about ride~1 of your predictions: a steady
$5\,\mathrm{m/s}$ upward is \emph{zero acceleration}, so the scale reads
exactly 60.0 --- the reading cares about acceleration, never about velocity.
If a sealed elevator cruises perfectly smoothly, no scale (and no sensation
in your stomach) can tell you whether you are moving at all. That is the
first law speaking.
\end{minipage}
\hfill
\begin{minipage}[t]{0.6\textwidth}
\vspace{0pt}
\centering
\begin{tikzpicture}[>=Stealth, scale=1.80]
    \fill[paper, rounded corners=5pt] (-1.55,-0.18) rectangle (1.55,3.76);
    \draw[very thick] (-1.3,0) rectangle (1.3,3.0);
    \draw[thick] (0,3.6) -- (0,3.0);
    \fill[gray!20] (-1.3,0) rectangle (1.3,0.10);
    \fill[gray!30] (-0.58,0.12) rectangle (0.58,0.40);
    \draw[thick]   (-0.58,0.12) rectangle (0.58,0.40);
    \fill[white]   (-0.32,0.15) rectangle (0.32,0.37);
    \node[font=\tiny] at (0,0.26) {60.0\,kg};
    \draw[thick] (-0.16,0.40) -- (0,0.62) -- (0.16,0.40);
    \draw[thick] (0,0.62) -- (0,1.52);
    \draw[thick] (0,1.05) -- (-0.42,0.85);
    \draw[thick] (0,1.05) -- ( 0.42,0.85);
    \draw[thick,fill=paper] (0,1.75) circle (0.23);
    \draw[very thick,-Stealth,myblue]  (-0.82,0.40) -- (-0.82,2.0);
    \node[left,myblue,font=\small] at (-0.87,1.20) {$\vec{\mathscr{N}}$};
    \draw[very thick,-Stealth,myred]   ( 0.80,2.00) -- ( 0.80,0.40);
    \node[right,myred,font=\small] at (0.85,1.20) {$m\vec{g}$};
\end{tikzpicture}
\label{fig:scale}
\end{minipage}

 
\vspace{2.4cm}
\noindent\textbf{\sffamily(b)\; Accelerating Upward.}
 
\noindent
\begin{minipage}[t]{0.55\textwidth}
\vspace{0pt}
Now the elevator accelerates upward. To make the person accelerate upward too, there must
be a net upward force on them. But the person's inertia resists this change in motion: their
body ``wants'' to stay at rest or moving as it was. The floor must therefore push up harder than before, hard enough both to support the weight and to supply the extra upward force the acceleration
demands. The normal force rises above $mg$. Newton's second law gives
\begin{equation}
    \mathscr{N} - mg = ma \implies \mathscr{N} = m(g+a).
\end{equation}
Suppose the scale reads $64.8\,\mathrm{kg}$, meaning $\mathscr{N} = 648\,\mathrm{N}$:
\begin{equation*}
    648 = 60(10+a) \implies a = 0.8\;\mathrm{m\,s^{-2}}\;\text{upward.}
\end{equation*}
The person feels heavier because the floor genuinely presses harder on their feet. The true
weight is unchanged at $600\,\mathrm{N}$; only the apparent weight has risen, to
$648\,\mathrm{N}$.

\end{minipage}%
\hfill
\begin{minipage}[t]{0.6\textwidth}
\vspace{0pt}
\centering
\begin{tikzpicture}[>=Stealth, scale=1.80]
    \fill[paper, rounded corners=5pt] (-1.55,-0.18) rectangle (2.15,3.7);
    \draw[very thick] (-1.3,0) rectangle (1.3,3.0);
    \draw[thick] (0,3.6) -- (0,3.0);
    \fill[gray!20] (-1.3,0) rectangle (1.3,0.10);
    \fill[gray!30] (-0.58,0.12) rectangle (0.58,0.40);
    \draw[thick]   (-0.58,0.12) rectangle (0.58,0.40);
    \fill[white]   (-0.32,0.15) rectangle (0.32,0.37);
    \node[font=\tiny] at (0,0.26) {64.8\,kg};
    \draw[thick] (-0.16,0.40) -- (0,0.62) -- (0.16,0.40);
    \draw[thick] (0,0.62) -- (0,1.52);
    \draw[thick] (0,1.05) -- (-0.42,0.85);
    \draw[thick] (0,1.05) -- ( 0.42,0.85);
    \draw[thick,fill=paper] (0,1.75) circle (0.23);
    \draw[very thick,-Stealth,myblue]  (-0.82,0.40) -- (-0.82,2.58);
    \node[left,myblue,font=\small] at (-0.87,1.49) {$\vec{\mathscr{N}}$};
    \draw[very thick,-Stealth,myred]   ( 0.80,2.00) -- ( 0.80,0.40);
    \node[right,myred,font=\small] at (0.85,1.20) {$m\vec{g}$};
    \draw[very thick,-Stealth,mygreen] (1.80,1.30) -- (1.80,2.90);
    \node[right,mygreen,font=\small] at (1.85,2.10) {$\vec{a}$};
\end{tikzpicture}
\end{minipage}



\vspace{0.3cm}
\noindent\textbf{\sffamily(c)\; Accelerating Downward.}
 
\noindent
\begin{minipage}[t]{0.55\textwidth}
\vspace{0pt}
When the elevator accelerates downward, the person's inertia resists the change: the body
"wants" to keep moving as it was, so for a moment it does not fall away with the floor as
readily as the floor drops beneath it. The person presses less firmly on the scale, and so
the scale presses less firmly back. The normal force falls below $mg$, leaving a net
downward force that produces the downward acceleration:
\begin{equation}
    mg - \mathscr{N} = ma \implies \mathscr{N} = m(g-a). \label{downward}
\end{equation}
If the scale reads $54.0\,\mathrm{kg}$, meaning $\mathscr{N} = 540\,\mathrm{N}$:
\begin{equation*}
    540 = 60(10-a) \implies a = 1.0\;\mathrm{m\,s^{-2}}\;\text{downward.}
\end{equation*}
The person feels lighter because the floor really does press less on their feet. The true
weight is unchanged at $600\,\mathrm{N}$; only the apparent weight has fallen, to
$540\,\mathrm{N}$.
\end{minipage}%
\hfill
\begin{minipage}[t]{0.60\textwidth}
\vspace{0pt}
\centering
\begin{tikzpicture}[>=Stealth, scale=1.8]
    \fill[paper, rounded corners=5pt] (-1.55,-0.18) rectangle (2.15,3.7);
    \draw[very thick] (-1.3,0) rectangle (1.3,3.0);
    \draw[thick] (0,3.6) -- (0,3.0);
    \fill[gray!20] (-1.3,0) rectangle (1.3,0.12);
    \fill[gray!30] (-0.58,0.12) rectangle (0.58,0.40);
    \draw[thick]   (-0.58,0.12) rectangle (0.58,0.40);
    \fill[white]   (-0.32,0.15) rectangle (0.32,0.37);
    \node[font=\tiny] at (0,0.26) {54.0\,kg};
    \draw[thick] (-0.16,0.40) -- (0,0.62) -- (0.16,0.40);
    \draw[thick] (0,0.62) -- (0,1.52);
    \draw[thick] (0,1.05) -- (-0.42,0.85);
    \draw[thick] (0,1.05) -- ( 0.42,0.85);
    \draw[thick,fill=paper] (0,1.75) circle (0.23);
    \draw[very thick,-Stealth,myblue]  (-0.82,0.40) -- (-0.82,1.60);
    \node[left,myblue,font=\small] at (-0.87,1.00) {$\vec{\mathscr{N}}$};
    \draw[very thick,-Stealth,myred]   ( 0.80,2.00) -- ( 0.80,0.40);
    \node[right,myred,font=\small] at (0.85,1.20) {$m\vec{g}$};
    \draw[very thick,-Stealth,mygreen] (1.80,2.90) -- (1.80,1.30);
    \node[right,mygreen,font=\small] at (1.85,2.10) {$\vec{a}$};
\end{tikzpicture}
\end{minipage}
 
\vspace{1.2cm}
\noindent\textbf{\sffamily(d)\; Free Fall.}
 
\noindent
\begin{minipage}[t]{0.55\textwidth}
\vspace{0pt}
Now imagine the cable snaps and the elevator falls freely. Gravity acts on everything
inside equally, so the person, the scale, and the floor all accelerate downward together
at the same rate $a = g$. Because the floor falls away just as fast as the person does, it
never has to push on them at all. The normal force drops all the way to zero. Setting
$a = g$ in equation~\eqref{downward},
\begin{equation*}
    \mathscr{N} = m(g - g) = 0.
\end{equation*}
The scale reads zero, yet gravity has certainly not switched off; the Earth still pulls the
person down with the full force $mg$. The person simply floats inside the falling car
because there is no longer any contact force to register. This is exactly what we mean by
\textit{weightlessness}: not the absence of gravity, but the absence of the normal force.
\end{minipage}%
\hfill
\begin{minipage}[t]{0.60\textwidth}
\vspace{0pt}
\centering
\begin{tikzpicture}[>=Stealth, scale=1.80]
    \fill[paper, rounded corners=5pt] (-1.55,0.22) rectangle (2.4,4.0);
    \draw[very thick] (-1.3,0.4) rectangle (1.3,3.4);
    \draw[thick] (0,3.4) -- (0,3.62);
    \draw[thick] (0,3.72) -- (0,3.92);
    \draw[thick, myred] (-0.10, 3.59) -- (0.12, 3.75);
    \fill[gray!20] (-1.3,0.4) rectangle (1.3,0.52);
    \fill[gray!30] (-0.58,0.68) rectangle (0.58,0.96);
    \draw[thick]   (-0.58,0.68) rectangle (0.58,0.96);
    \fill[white]   (-0.32,0.71) rectangle (0.32,0.93);
    \node[font=\tiny] at (0,0.82) {0.0\,kg};
    \draw[dashed,gray!55,thin] (-0.58,0.52) -- (-0.58,0.68);
    \draw[dashed,gray!55,thin] ( 0.58,0.52) -- ( 0.58,0.68);
    \draw[thick] (-0.16,1.28) -- (0,1.50) -- (0.16,1.28);
    \draw[thick] (0,1.50) -- (0,2.40);
    \draw[thick] (0,1.93) -- (-0.42,1.73);
    \draw[thick] (0,1.93) -- ( 0.42,1.73);
    \draw[thick,fill=paper] (0,2.63) circle (0.23);
    \draw[dashed,gray!55,thin] (-0.16,0.96) -- (-0.16,1.28);
    \draw[dashed,gray!55,thin] ( 0.16,0.96) -- ( 0.16,1.28);
    \draw[very thick,-Stealth,myred] (0.80,2.83) -- (0.80,1.21);
    \node[right,myred,font=\small] at (0.85,2.02) {$m\vec{g}$};
    \draw[very thick,-Stealth,mygreen] (1.80,3.20) -- (1.80,1.60);
    \node[right,mygreen,font=\small] at (1.85,2.40) {$\vec{a}{=}\vec{g}$};
\end{tikzpicture}
\end{minipage}
 


\begin{checkpoint}[5.1: The floating astronaut]
Astronauts on the International Space Station float freely inside the cabin.
The station orbits about 400\,km above the Earth's surface. The astronauts
float because
\begin{enumerate}[label=(\alph*), leftmargin=2.5em, itemsep=0pt]
  \item at that altitude they are beyond the reach of the Earth's gravity.
  \item the Earth's gravity is cancelled by the pull of the Moon and Sun.
  \item they and the station are both in free fall, falling around the Earth together.
  \item the station's speed generates an outward force that balances gravity.
\end{enumerate}
\checkanswer{Answer: (c). At 400 km up, the Earth's gravity is still about
90\% of its surface strength --- answer (a) is the single most common
misconception in this course. The station and everyone in it are permanently
in case (d) of the elevator: falling freely, together, all the time. Their
sideways speed simply makes them keep missing the Earth --- which is what an
orbit is, as the section on gravitation will make precise.}
\end{checkpoint}

\begin{supplementary}{A Glimpse of Something Deeper: The Equivalence Principle}

We close the elevator story with a question that troubled Einstein and ultimately led him to general relativity.
Imagine you wake up inside a sealed box with no windows. You feel a normal force pressing
up on your feet. How could you tell whether you are sitting stationary on the surface of the
Earth, with gravity pulling you down, or drifting in empty space far from any planet, with
the box accelerating upward at $g=9.8\,\text{m/s}^2$?

The remarkable answer is that you cannot. No mechanical experiment conducted inside the
box will distinguish between these two situations. Drop a ball: in both cases it accelerates
toward the floor at $10\,\text{m/s}^2$. Step on a scale: it reads $mg$ in both cases. Einstein
elevated this to a fundamental principle: \textbf{gravity and accelerated motion are locally
equivalent}. This is the \textbf{equivalence principle}.

\textbf{a) Two Kinds of Mass.} The symbol $m$ appears in two logically distinct roles.
In Newton's second law, $m$ appears as \textbf{inertial mass}: the resistance of an object to
acceleration. In the gravitational force law, $m$ appears as \textbf{gravitational mass}: the
quantity that determines how strongly gravity pulls on the object. There is no obvious reason
these two quantities should be equal. Yet every experiment ever performed finds them equal.
The Hungarian physicist Lor\'and E\"otv\"os demonstrated this equality to one part in $10^8$
in the 1880s, and modern experiments push this to one part in $10^{15}$.

The equality $m_\text{inertial}=m_\text{gravitational}$ is exactly what ensures that all
objects fall at the same rate regardless of their mass.

\textbf{b) Light Must Bend.} The equivalence principle has a striking consequence that
goes beyond mechanics. Shine a horizontal beam of light across the accelerating box. While
the light travels across, the box accelerates upward, so the beam hits the far wall slightly
below where it was aimed --- the light beam appears to curve downward. Since an accelerating
box is locally indistinguishable from a gravitational field, gravity must also bend light. This
was a radical prediction: Newtonian gravity acts only on objects with mass, and light was
thought to be massless. Einstein's prediction was confirmed experimentally in 1919 by Arthur
Eddington, who photographed stars near the edge of the Sun during a solar eclipse and found
that their apparent positions were shifted by exactly the amount Einstein predicted. The
announcement made Einstein famous overnight.

\textbf{c) Where This Leads.} Once you take the equivalence principle seriously and ask what
geometry of space and time is consistent with it, you are led to \textbf{general relativity}.
Gravity is reinterpreted not as a force but as the curvature of spacetime caused by mass and
energy. Objects in free fall follow the straightest possible paths through this curved geometry.
The bending of light, the slowing of clocks in gravitational fields, the existence of
gravitational waves, and the physics of black holes all follow from the same geometric
framework. The elevator we analysed today is the first step in that journey.
\end{supplementary}







\newpage
\section{Circular Motion}

We shall now investigate a special class of motions: motion in a plane about a central point, which we refer to as \textit{central motion}. The most important special case is circular motion. There are many instances of circular motion: a bicycle rider on a circular track, a ball spun around on a string, and the rotation of a spinning wheel are just a few examples. The motion of the Moon around the Earth is nearly circular. The motions of the planets around the Sun are nearly circular. Our Sun moves in a nearly circular orbit about the centre of our galaxy.

We begin by describing the kinematics of circular motion -- the position, velocity, and acceleration -- as a special case of two-dimensional motion. Unlike linear motion, where velocity and acceleration are directed along the line of motion, in circular motion the velocity is always tangent to the circle. This means that as the object moves in a circle, the direction of the velocity is always changing. Examining this change reveals that the direction of the acceleration is towards the centre of the circle. This inward component of acceleration is called the \textit{centripetal acceleration}. If the object is also changing speed, there is an additional non-zero tangential acceleration in the direction of motion. For uniform circular motion (constant speed), only the centripetal acceleration is non-zero.

Newton recognised that the Moon is always ``falling'' towards the Earth; otherwise, by the First Law, it would continue in a straight line rather than follow a circular orbit. There must therefore exist a centripetal force, a radial force pointing inward, producing the centripetal acceleration. Because Newton's Second Law is a vector equation, the radial component gives
\begin{equation}
    F_r = m a_r.
\end{equation}

% ==============================================================
\subsection{Circular Motion: Velocity and Angular Velocity}
% ==============================================================

We describe circular motion using polar coordinates. The position vector $\vec{r}(t)$ of an object moving in a circular orbit of radius $r$ is written as
\begin{equation}
    \vec{r}(t) = r\,\hat{r}(t),
\end{equation}
where $\hat{r}(t)$ is the radial unit vector pointing from the origin to the object. At any point $P$ with coordinates $(r, \theta(t))$, we define two sets of unit vectors: the polar basis $(\hat{r}(t),\, \hat{\theta}(t))$ and the Cartesian basis $(\hat{\imath},\, \hat{\jmath})$, as shown in Figure~\ref{fig:circleunitvec}. Their relationship is
\begin{align}
    \hat{r}(t)      &= \cos\theta(t)\,\hat{\imath} + \sin\theta(t)\,\hat{\jmath}, \label{eq:rhat}\\
    \hat{\theta}(t) &= -\sin\theta(t)\,\hat{\imath} + \cos\theta(t)\,\hat{\jmath}. \label{eq:thetahat}
\end{align}

\begin{figure}[ht]
\centering
\begin{tikzpicture}[>=Stealth, thick, scale=1.15]
    % Axes
    \draw[->, thin, gray!70] (-2.6,0) -- (2.8,0) node[right, black, font=\small] {$+x$};
    \draw[->, thin, gray!70] (0,-2.6) -- (0,2.8) node[above, black, font=\small] {$+y$};
    % Cartesian unit vector labels
    \node[font=\scriptsize, gray!60!black] at (0.55, 0.18) {$\hat{\imath}$};
    \node[font=\scriptsize, gray!60!black] at (0.18, 0.55) {$\hat{\jmath}$};
    % Circle
    \draw[dashed, gray!55] (0,0) circle (2);
    % Origin
    \fill (0,0) circle (1.5pt);
    \node[below left, font=\small] at (0,0) {$O$};
    % Point P
    \coordinate (P) at (52:2);
    \fill[black] (P) circle (2pt);
    \node[above right, font=\small] at (P) {$P$};
    % Position vector
    \draw[->, very thick] (0,0) -- (P) node[midway, left, font=\small] {$\vec{r}(t)$};
    % Angle arc and label
    \draw[thin] (0.65,0) arc (0:52:0.65);
    \node[font=\scriptsize] at (26:0.92) {$\theta(t)$};
    % Radius label
    \node[font=\scriptsize, gray!50!black] at (74:2.22) {$r$};
    % r-hat unit vector
    \draw[->, thick] (P) -- ++(52:0.80) node[above right, font=\small] {$\hat{r}(t)$};
    % theta-hat unit vector (CCW tangent)
    \draw[->, thick] (P) -- ++(142:0.80) node[above left, font=\small] {$\hat{\theta}(t)$};
\end{tikzpicture}
\caption{A circular orbit of radius $r$ with the polar unit vectors $\hat{r}(t)$ and $\hat{\theta}(t)$ at the point $P$.}
\label{fig:circleunitvec}
\end{figure}

Before computing the velocity, we calculate the time derivatives of the unit vectors. Using the chain rule on Eq.~(\ref{eq:rhat}),
\begin{equation}
    \frac{d\hat{r}(t)}{dt} = \left(-\sin\theta(t)\,\hat{\imath} + \cos\theta(t)\,\hat{\jmath}\right)\frac{d\theta}{dt} = \frac{d\theta}{dt}\,\hat{\theta}(t).
\end{equation}
Similarly, differentiating Eq.~(\ref{eq:thetahat}),
\begin{equation}
    \frac{d\hat{\theta}(t)}{dt} = \left(-\cos\theta(t)\,\hat{\imath} - \sin\theta(t)\,\hat{\jmath}\right)\frac{d\theta}{dt} = -\frac{d\theta}{dt}\,\hat{r}(t).
\end{equation}
The velocity vector is then
\begin{equation}
    \vec{v}(t) = \frac{d\vec{r}}{dt} = r\frac{d\hat{r}}{dt} = r\frac{d\theta}{dt}\,\hat{\theta}(t) = v_\theta\,\hat{\theta}(t),
\end{equation}
where the tangential component of the velocity is
\begin{equation}
    v_\theta = r\frac{d\theta}{dt}.
\end{equation}
The \textit{angular speed} is defined as the magnitude of the rate of change of angle,
\begin{equation}
    \omega \equiv \left|\frac{d\theta}{dt}\right|,
\end{equation}
so that the speed is $v = |\vec{v}| = r\omega$.

% --------------------------------------------------------------
\subsection{Geometric Derivation of the Velocity}
% --------------------------------------------------------------

Consider a particle undergoing circular motion. During a time interval $\Delta t$, the particle moves from position $\vec{r}(t)$ to $\vec{r}(t+\Delta t)$, sweeping an angle $\Delta\theta$, as shown in Figure~\ref{fig:displacement}.

\begin{figure}[ht]
\centering
\begin{tikzpicture}[>=Stealth, thick, scale=1.2]
    % Arc
    \draw[dashed, gray!55] (0,0) circle (2);
    \fill (0,0) circle (1.5pt) node[below left, font=\small] {$O$};
    % Two position vectors
    \coordinate (A) at (20:2);
    \coordinate (B) at (65:2);
    \draw[->, thick] (0,0) -- (A) node[right, font=\small] {$\vec{r}(t)$};
    \draw[->, thick] (0,0) -- (B) node[above right, font=\small] {$\vec{r}(t{+}\Delta t)$};
    \fill (A) circle (1.5pt);
    \fill (B) circle (1.5pt);
    % Displacement chord
    \draw[->, thick, red] (A) -- (B) node[midway, right, font=\small] {$\Delta\vec{r}$};
    % Angle arc
    \draw[thin] (0.7,0) arc (0:20:0.7);
    \draw[thin] (0.55,0) arc (0:65:0.55);
    \node[font=\scriptsize] at (42:0.82) {$\Delta\theta$};
\end{tikzpicture}
\caption{Displacement vector $\Delta\vec{r}$ for circular motion over an interval $\Delta t$.}
\label{fig:displacement}
\end{figure}

The magnitude of the displacement is $|\Delta\vec{r}| = 2r\sin(\Delta\theta/2)$. For small angles, $\sin(\Delta\theta/2) \approx \Delta\theta/2$, so
\begin{equation}
    |\Delta\vec{r}| \approx r\,\Delta\theta.
\end{equation}
This is consistent with the limit $\Delta\theta \to 0$, where the chord approaches the arc length $r\,\Delta\theta$. The magnitude of the velocity then follows as
\begin{equation}
    v \equiv |\vec{v}(t)| = \lim_{\Delta t\to 0}\frac{|\Delta\vec{r}|}{\Delta t} = r\lim_{\Delta t\to 0}\frac{\Delta\theta}{\Delta t} = r\frac{d\theta}{dt} = r\omega.
\end{equation}
As $\Delta t\to 0$, the direction of $\Delta\vec{r}$ approaches the tangent to the circle, perpendicular to $\vec{r}(t)$ and in the $+\hat{\theta}$-direction for counterclockwise motion. This confirms that the velocity is always tangent to the circle.

% ==============================================================
\subsection{Tangential and Radial Acceleration}
% ==============================================================

In polar coordinates, the acceleration has two components: a tangential component $a_\theta$ and a radial component $a_r$,
\begin{equation}
    \vec{a}(t) = a_r\,\hat{r}(t) + a_\theta\,\hat{\theta}(t).
\end{equation}
The unit vectors $\hat{r}(t)$ and $\hat{\theta}(t)$ change direction with time and are therefore not constant. To compute the acceleration, we differentiate the velocity $\vec{v}(t) = r(d\theta/dt)\,\hat{\theta}(t)$ using the product rule,
\begin{equation}
    \vec{a}(t) = r\frac{d^2\theta}{dt^2}\,\hat{\theta}(t) + r\frac{d\theta}{dt}\frac{d\hat{\theta}}{dt}.
\end{equation}
Substituting $d\hat{\theta}/dt = -(d\theta/dt)\hat{r}(t)$,
\begin{equation}
    \vec{a}(t) = r\frac{d^2\theta}{dt^2}\,\hat{\theta}(t) - r\left(\frac{d\theta}{dt}\right)^2\hat{r}(t).
\end{equation}
The two components are therefore
\begin{align}
    a_\theta &= r\frac{d^2\theta}{dt^2}, \label{eq:atangential}\\
    a_r      &= -r\left(\frac{d\theta}{dt}\right)^2 = -r\omega^2 < 0. \label{eq:aradial}
\end{align}
Because $a_r < 0$, the radial component $a_r(t) = -r\omega^2\,\hat{r}(t)$ is always directed towards the centre of the circular orbit. This is the centripetal acceleration.

\begin{workedexample}[5.2: Circular Motion Kinematics]
A particle moves in a circle of radius $R$. At $t = 0$ it is on the $x$-axis. The angle with the positive $x$-axis is $\theta(t) = At^3 - Bt$, where $A$ and $B$ are positive constants. Find (a) the velocity vector and (b) the acceleration vector in polar coordinates. At what time is the centripetal acceleration zero?

\noindent\textbf{Solution.}
The derivatives are $d\theta/dt = 3At^2 - B$ and $d^2\theta/dt^2 = 6At$. The velocity vector is
\[
    \vec{v}(t) = R\frac{d\theta}{dt}\,\hat{\theta}(t) = R(3At^2 - B)\,\hat{\theta}(t).
\]
The acceleration vector is
\[
    \vec{a}(t) = R\frac{d^2\theta}{dt^2}\,\hat{\theta}(t) - R\left(\frac{d\theta}{dt}\right)^2\hat{r}(t) = 6RAt\,\hat{\theta}(t) - R(3At^2 - B)^2\hat{r}(t).
\]
The centripetal acceleration is zero when $d\theta/dt = 0$, that is, when
\[
    3At_1^2 - B = 0 \implies t_1 = \sqrt{\frac{B}{3A}}.
\]
\end{workedexample}

% ==============================================================
\subsection{Period and Frequency for Uniform Circular Motion}
% ==============================================================

If the total tangential force on the object is zero, the tangential acceleration vanishes,
\begin{equation}
    a_\theta = 0,
\end{equation}
and the speed remains constant. This is \textit{uniform circular motion}. The acceleration reduces to
\begin{equation}
    \vec{a}_r(t) = -r\omega^2\,\hat{r}(t) \qquad \text{(uniform circular motion)}.
\end{equation}
Because the speed $v = r\omega$ is constant, the time to complete one full revolution, the \textit{period} $T$, is constant. In one period the object travels the circumference $2\pi r$, so
\begin{equation}
    2\pi r = vT \implies T = \frac{2\pi r}{v} = \frac{2\pi r}{r\omega} = \frac{2\pi}{\omega}.
\end{equation}
The \textit{frequency} $f$ is the reciprocal of the period,
\begin{equation}
    f = \frac{1}{T} = \frac{\omega}{2\pi}.
\end{equation}
The SI unit of frequency is the hertz: $[\mathrm{Hz}] = [\mathrm{s}^{-1}]$. The magnitude of the centripetal acceleration can be written in several equivalent forms. Using $v = r\omega$,
\begin{align}
    |a_r| &= r\omega^2, \label{eq:ar1}\\
    |a_r| &= \frac{v^2}{r}, \label{eq:ar2}\\
    |a_r| &= 4\pi^2 r f^2, \label{eq:ar3}\\
    |a_r| &= \frac{4\pi^2 r}{T^2}. \label{eq:ar4}
\end{align}
Which form to use depends on what information is given about the orbit.

% --------------------------------------------------------------
\subsection{Geometric Interpretation of the Radial Acceleration}
% --------------------------------------------------------------

An object in circular orbit is always accelerating towards the centre. The velocity is always tangent to the circle, so its direction is continuously changing as the object moves. Because the velocity changes direction, the acceleration is non-zero even when the speed is constant.

\begin{figure}[ht]
\centering
\begin{tikzpicture}[>=Stealth, thick, scale=1.1]
    % Left: circle with two velocity vectors
    \begin{scope}[shift={(-3.5,0)}]
        \draw[dashed, gray!55] (0,0) circle (1.8);
        \fill (0,0) circle (1.5pt);
        % Point at angle 30
        \coordinate (A) at (30:1.8);
        \coordinate (B) at (80:1.8);
        \fill (A) circle (2pt);
        \fill (B) circle (2pt);
        % Position vectors
        \draw[->, thin, gray!60] (0,0) -- (A);
        \draw[->, thin, gray!60] (0,0) -- (B);
        % Velocity at A (tangent, 90 deg CCW)
        \draw[->, thick] (A) -- ++(120:1.0) node[right, font=\scriptsize] {$\vec{v}(t)$};
        % Velocity at B
        \draw[->, thick] (B) -- ++(170:1.0) node[left, font=\scriptsize] {$\vec{v}(t{+}\Delta t)$};
        % Angle label
        \draw[thin] (0.55,0) arc (0:30:0.55);
        \draw[thin] (0.42,0) arc (0:80:0.42);
        \node[font=\scriptsize] at (55:0.75) {$\Delta\theta$};
        % Labels for position vectors
        \node[right, font=\scriptsize] at ($(0,0)!0.55!(A)$) {$\vec{r}(t)$};
        \node[above left, font=\scriptsize] at ($(0,0)!0.55!(B)$) {$\vec{r}(t{+}\Delta t)$};
    \end{scope}
    % Right: velocity triangle
    \begin{scope}[shift={(3.0,0)}]
        % v(t) pointing up
        \coordinate (O2) at (0,0);
        \draw[->, thick] (O2) -- ++(90:1.6) node[right, font=\scriptsize] {$\vec{v}(t)$};
        % v(t+dt) at angle Delta theta from v(t): pointing up-left
        \draw[->, thick] (O2) -- ++(140:1.6) node[left, font=\scriptsize] {$\vec{v}(t{+}\Delta t)$};
        % Delta v (from tip of v(t+dt) to tip of v(t))
        \coordinate (Vtip) at ($(O2)+(90:1.6)$);
        \coordinate (Vdttip) at ($(O2)+(140:1.6)$);
        \draw[->, thick, red] (Vdttip) -- (Vtip) node[right, font=\scriptsize] {$\Delta\vec{v}$};
        % Angle label
        \draw[thin] (0,0.55) arc (90:140:0.55);
        \node[font=\scriptsize] at (115:0.75) {$\Delta\theta$};
    \end{scope}
\end{tikzpicture}
\caption{Left: velocity vectors at two nearby points on a circular orbit. Right: the change in velocity $\Delta\vec{v}$, which points inward as $\Delta\theta\to 0$.}
\label{fig:velchange}
\end{figure}

The magnitude of the change in velocity is $|\Delta\vec{v}| = 2v\sin(\Delta\theta/2) \approx v\,\Delta\theta$ for small $\Delta\theta$ (Figure~\ref{fig:velchange}). The magnitude of the radial acceleration is therefore
\begin{equation}
    |a_r| = \lim_{\Delta t\to 0}\frac{|\Delta\vec{v}|}{\Delta t} = v\lim_{\Delta t\to 0}\frac{\Delta\theta}{\Delta t} = v\frac{d\theta}{dt} = v\omega = \frac{v^2}{r}.
\end{equation}
As $\Delta\theta\to 0$, the direction of $\Delta\vec{v}$ is perpendicular to $\vec{v}$ and directed inward, in the $-\hat{r}$-direction, confirming the centripetal character of the acceleration.

% ==============================================================
\subsection{Angular Velocity and Angular Acceleration}
% ==============================================================

% --------------------------------------------------------------
\subsection{Angular Velocity}
% --------------------------------------------------------------

We always choose a right-handed cylindrical coordinate system. With the positive $z$-axis pointing up, the angle $\theta$ increases counterclockwise as viewed from above. The \textit{angular velocity vector} for a point object undergoing circular motion about the $z$-axis is directed along the $z$-axis,
\begin{equation}
    \vec{\omega} = \frac{d\theta}{dt}\,\hat{k} = \omega_z\,\hat{k}.
\end{equation}
The angular speed is the magnitude of the $z$-component, $\omega \equiv |\omega_z| = |d\theta/dt|$. If the object rotates counterclockwise, $\omega_z = d\theta/dt > 0$ and $\vec{\omega}$ points in the $+\hat{k}$-direction. If the object rotates clockwise, $\omega_z < 0$ and $\vec{\omega}$ points in the $-\hat{k}$-direction.

The velocity and angular velocity are related by the cross product,
\begin{equation}
    \vec{v} = \vec{\omega}\times\vec{r} = \frac{d\theta}{dt}\,\hat{k}\times r\hat{r} = r\frac{d\theta}{dt}\,\hat{\theta}.
\end{equation}

\begin{workedexample}[5.3: Angular Velocity]
A particle moves in a circle of radius $R$. At $t = 0$ it is on the $x$-axis. The angle with the positive $x$-axis is $\theta(t) = At - Bt^3$, where $A$ and $B$ are positive constants. Find (a) the angular velocity vector, (b) the velocity vector, (c) the time $t_1$ when the angular velocity is zero, and (d) the direction of $\vec{\omega}$ for $t < t_1$ and $t > t_1$.

\noindent\textbf{Solution.} The derivative of $\theta(t)$ is
\[
    \frac{d\theta}{dt} = A - 3Bt^2.
\]
(a) The angular velocity vector is
\[
    \vec{\omega}(t) = (A - 3Bt^2)\,\hat{k}.
\]
(b) The velocity is
\[
    \vec{v}(t) = R(A - 3Bt^2)\,\hat{\theta}(t).
\]
(c) The angular velocity is zero when
\[
    A - 3Bt_1^2 = 0 \implies t_1 = \sqrt{\frac{A}{3B}}.
\]
(d) For $t < t_1$, $d\theta/dt > 0$ so $\vec{\omega}$ points in the $+\hat{k}$-direction (counterclockwise). For $t > t_1$, $d\theta/dt < 0$ so $\vec{\omega}$ points in the $-\hat{k}$-direction (clockwise).
\end{workedexample}

% --------------------------------------------------------------
\subsection{Angular Acceleration}
% --------------------------------------------------------------

The \textit{angular acceleration vector} for circular motion about the $z$-axis is
\begin{equation}
    \vec{\alpha} = \frac{d^2\theta}{dt^2}\,\hat{k} = \alpha_z\,\hat{k}.
\end{equation}
The SI units of angular acceleration are $[\mathrm{rad\cdot s^{-2}}]$. There are four combinations of the signs of $\omega_z$ and $\alpha_z$ to consider.

When $\vec{\alpha}$ points in the $+\hat{k}$-direction ($\alpha_z > 0$): (i) if $\omega_z > 0$ (counterclockwise), the object is speeding up counterclockwise; (ii) if $\omega_z < 0$ (clockwise), the object is slowing down.

When $\vec{\alpha}$ points in the $-\hat{k}$-direction ($\alpha_z < 0$): (iii) if $\omega_z > 0$ (counterclockwise), the object is slowing down; (iv) if $\omega_z < 0$ (clockwise), the object is speeding up clockwise.

The tangential and radial components of the acceleration (Eqs.~(\ref{eq:atangential}) and (\ref{eq:aradial})) can be written compactly as $a_\theta = r\alpha_z$ and $a_r = -r\omega_z^2$.

\begin{workedexample}[5.4: Integration and Circular Motion Kinematics]
A point object is constrained to travel in a circle. The $z$-component of the angular acceleration is
\[
    \alpha_z(t) = \begin{cases} b\!\left(1 - \dfrac{t}{t_1}\right), & 0 \leq t \leq t_1, \\[4pt] 0, & t > t_1, \end{cases}
\]
where $b > 0$ has units $\mathrm{rad\cdot s^{-2}}$. (a) Find the angular velocity at $t = t_1$. (b) Through what angle has the object rotated at $t = t_1$?

\noindent\textbf{Solution.} (a) Starting from rest ($\omega_z(0) = 0$), we integrate:
\[
    \omega_z(t_1) = \int_0^{t_1} b\!\left(1 - \frac{t'}{t_1}\right)dt' = b\!\left[t_1 - \frac{t_1}{2}\right] = \frac{bt_1}{2}.
\]

(b) First find $\omega_z(t)$ for $0 \leq t \leq t_1$:
\[
    \omega_z(t) = \int_0^t b\!\left(1 - \frac{t'}{t_1}\right)dt' = b\!\left(t - \frac{t^2}{2t_1}\right).
\]
Then integrate again to find the angle:
\[
    \theta(t_1) - \theta(0) = \int_0^{t_1}\omega_z(t')\,dt' = \int_0^{t_1} b\!\left(t' - \frac{t'^2}{2t_1}\right)dt' = b\!\left[\frac{t_1^2}{2} - \frac{t_1^2}{6}\right] = \frac{bt_1^2}{3}.
\]
\end{workedexample}


\newpage
% ==============================================================
\section{Gravitation: The Moon is Falling}
% ==============================================================

At the start of our circular-motion story we repeated Newton's remark: the
Moon is always \emph{falling} towards the Earth --- otherwise, by the first
law, it would fly off in a straight line. We now have every tool needed to
take that remark seriously, and it leads to one of the greatest
predict-and-check episodes in the history of science. This section retells
it, and asks you to redo Newton's arithmetic yourself.

In 1665 the plague closed Cambridge University, and the 23-year-old Newton
retreated to his family farm at Woolsthorpe. By his own (much later) telling,
watching an apple fall in the garden set him wondering: the Earth's pull
evidently reaches the top of the tallest tree --- might it reach all the way
to the Moon? If so, the Moon's endless circling and the apple's brief drop
would be \emph{the same phenomenon}. No one had ever connected the heavens
and the ground by a single law. The idea is bold; what makes it science is
that it can be checked by a number.

\begin{predictbox}[5.3: Newton's own prediction]
Suppose the Earth's pull weakens with the \emph{square} of the distance from
the Earth's centre --- Newton's guess, partly inspired by how the light of a
lamp dilutes with distance. An apple sits about one Earth radius from the
centre ($R_E \approx 6.37\times10^6$\,m). The Moon sits about
\emph{sixty} Earth radii away. Under the inverse-square guess:
\begin{enumerate}[leftmargin=2em, itemsep=1pt]
  \item How many times weaker should gravity be at the Moon than at the apple?
  \item The apple accelerates at $g = 9.8\;\mathrm{m/s^2}$. Predict the Moon's
        acceleration.
\end{enumerate}
Answer both before reading on --- then we will \emph{measure} the Moon's actual
acceleration, using nothing but the circular-motion formulas you just learned
and two numbers astronomers have known since antiquity.
\end{predictbox}

\subsection{The Moon Test}

If the inverse-square guess is right, gravity at the Moon should be weaker by
a factor of $60^2 = 3600$, so the Moon should be falling with acceleration
\begin{equation*}
    a_{\text{predicted}} = \frac{g}{3600} = \frac{9.8}{3600}
    \approx 2.7\times10^{-3}\;\mathrm{m/s^2}.
\end{equation*}
That is the prediction. Now the check. The Moon moves in a (nearly) circular
orbit, and for circular motion we know exactly what the centripetal
acceleration must be --- Eq.~\eqref{eq:ar4}:
\begin{equation*}
    |a_r| = \frac{4\pi^2 r}{T^2}.
\end{equation*}
The orbital radius is $r = 3.84\times10^8$\,m and the orbital period is
$T = 27.3\;\text{days} = 2.36\times10^6$\,s. Therefore the Moon's actual
acceleration is
\begin{equation*}
    |a_r| = \frac{4\pi^2 (3.84\times10^8)}{(2.36\times10^6)^2}
    \approx 2.7\times10^{-3}\;\mathrm{m/s^2}.
\end{equation*}
The prediction and the measurement agree. Newton, recalling the calculation
years later, wrote that he ``compared the force requisite to keep the Moon in
her orb with the force of gravity at the surface of the earth, and found them
answer pretty nearly.'' \emph{Pretty nearly} --- the understatement of the
seventeenth century. The apple and the Moon fall by the same law, and the
law is inverse-square.

Notice what kind of argument this is: guess a law, extract a numerical
prediction, and confront it with a measurement that had no right to agree by
accident. A factor of 3600, predicted from a hunch about lamplight, confirmed
to within a percent. You will use exactly this pattern --- predict, then
check --- in your own experiments this year, if with slightly smaller stakes.

\subsection{The Law of Universal Gravitation}

Newton generalised the moon test into a law that applies to every pair of
masses in the universe --- two galaxies, or two sesame seeds.

\begin{keybox}[Newton's Law of Universal Gravitation]
Any two point masses $m_1$ and $m_2$ separated by a distance $r$ attract each
other along the line joining them with a force of magnitude
\begin{equation}
    F_g = G\,\frac{m_1 m_2}{r^2},
\end{equation}
where $G$ is the \textbf{universal gravitational constant},
\begin{equation*}
    G = 6.674\times10^{-11}\;\mathrm{N\,m^2\,kg^{-2}}.
\end{equation*}
By the Third Law the two masses feel equal and opposite pulls: the Earth
pulls the apple, and the apple pulls the Earth, with the same force.
\end{keybox}

Two remarks. First, for a sphere (like the Earth) the law applies as if all
the mass were concentrated at the centre --- a theorem that cost Newton real
effort to prove, and part of why he delayed publishing for twenty years.
That is why the apple's distance counts as one Earth \emph{radius}, measured
from the centre, not from the ground. Second, the constant $G$ is
astonishingly small: gravity is by far the weakest force you will meet in
this course. Two 1\,kg masses held 1\,m apart attract with about
$10^{-10}$\,N --- roughly the weight of a single flake of dust. Gravity runs
the universe only because masses like planets are enormous, and because,
unlike electricity, gravity has no negative charge to cancel it.

The value of $G$ was not measured in Newton's lifetime. In 1798 Henry
Cavendish --- a man so shy he communicated with his housekeeper by notes ---
suspended a dumbbell of small lead spheres from a fine wire and measured the
twist produced by the gravitational pull of two larger spheres, a whisper of
a force in a sealed room. Knowing $G$, one can put the Earth itself on the
balance: from $g$ and $R_E$, the Earth's mass follows at once (try the next
example). Cavendish's contemporaries called his experiment ``weighing the
Earth.''

\begin{workedexample}[5.5: Weighing the Earth]
Given $g = 9.81\;\mathrm{m/s^2}$ at the surface, $R_E = 6.37\times10^6$\,m,
and Cavendish's $G$, find the mass of the Earth.

\medskip
\textit{Solution.} A mass $m$ at the surface feels the attraction of the
whole Earth acting from its centre:
\begin{equation*}
    mg = G\,\frac{M_E\,m}{R_E^2}
    \implies M_E = \frac{g R_E^2}{G}
    = \frac{(9.81)(6.37\times10^6)^2}{6.674\times10^{-11}}
    \approx 5.97\times10^{24}\;\mathrm{kg}.
\end{equation*}
Notice that $m$ cancelled: \emph{every} mass at the surface accelerates at
the same $g = GM_E/R_E^2$, which is the deep reason (first noticed by
Galileo) that heavy and light objects fall together. Notice also what the
formula says about altitude: at the International Space Station,
$r = R_E + 400\,\mathrm{km}$, and $g$ drops only to about
$8.7\;\mathrm{m/s^2}$ --- 89\% of its surface value, exactly as Checkpoint
5.1 claimed.
\end{workedexample}

\subsection{Circular Orbits}

For a satellite of mass $m$ in a circular orbit of radius $r$ about the Earth,
gravity is the only force, and it must supply exactly the centripetal
acceleration $v^2/r$ that circular motion demands. Newton's second law in the
radial direction gives
\begin{equation}
    G\,\frac{M_E\,m}{r^2} = \frac{m v^2}{r}
    \implies v = \sqrt{\frac{G M_E}{r}}.
\end{equation}
The satellite's mass cancels again --- a heavy satellite and a paint flake
share any given orbit at the same speed. Writing $v = 2\pi r/T$ and squaring:
\begin{equation}
    T^2 = \frac{4\pi^2}{G M_E}\,r^3. \label{eq:kepler3}
\end{equation}
The square of the period is proportional to the \emph{cube} of the orbital
radius. Keep this result in view; it is about to connect us to an astronomer
who died twelve years before Newton was born.

\subsection{Kepler's Laws}

Between 1576 and 1601, on an island observatory financed by the Danish crown,
Tycho Brahe measured the positions of the planets night after night for
twenty-five years --- the finest naked-eye data ever taken, accurate to about
one-sixtieth of a degree. Johannes Kepler, his difficult and brilliant
assistant, inherited the data and spent nearly a decade in ``war with Mars,''
trying orbit after orbit until the numbers finally surrendered. What they
surrendered to are three laws, published in 1609 and 1619. We state them here;
deriving them fully belongs to a later chapter, but you should know what they
say and where they came from.

\begin{keybox}[Kepler's Laws of Planetary Motion]
\begin{enumerate}[leftmargin=2em, itemsep=2pt]
  \item \textbf{The Law of Ellipses.} Each planet moves on an \emph{ellipse}
        with the Sun at one focus. (A circle is the special case of an
        ellipse with both foci together.)
  \item \textbf{The Law of Equal Areas.} The line from the Sun to a planet
        sweeps out equal areas in equal times --- so a planet moves fastest
        when nearest the Sun and slowest when farthest.
  \item \textbf{The Harmonic Law.} The square of a planet's orbital period is
        proportional to the cube of its mean distance from the Sun:
        $T^2 \propto a^3$, with the \emph{same} constant of proportionality
        for every planet.
\end{enumerate}
\end{keybox}

Kepler found these laws empirically --- pure pattern, no mechanism. He never
knew \emph{why} they were true, though he suspected the Sun. Now look back at
Eq.~\eqref{eq:kepler3}: for circular orbits, the third law fell out of the
law of gravitation in two lines of algebra, complete with the constant
$4\pi^2/GM$ --- the same for every body orbiting the same central mass, just
as Kepler said. (Done with ellipses instead of circles, the same follows
exactly; that computation is beyond us this year.) This was the argument that
convinced Newton's contemporaries: one law, guessed from a falling apple,
reproduced three laws distilled from twenty-five years of Tycho's data. The
second law, the equal-areas rule, has an even prettier origin --- it turns out
to be a disguised conservation law, and it will fall into our hands almost
for free when we meet angular momentum in the next chapter.

\begin{whiteboard}[5.1: The satellite that stands still]
Television dishes in Bhutan point at a fixed spot in the sky. The satellite
they aim at is \emph{geostationary}: it hangs over the same point on the
equator forever.
\begin{enumerate}[label=(\alph*), leftmargin=2.5em, itemsep=1pt]
  \item Before computing: what must such a satellite's orbital period be?
  \item Use Eq.~\eqref{eq:kepler3} to find the orbital radius, and hence the
        altitude above the ground. ($M_E = 5.97\times10^{24}$\,kg; the
        sidereal day is $T = 86\,164$\,s.)
  \item A minister asks for a geostationary satellite parked directly above
        Paro (latitude $27^\circ$N), to improve coverage. As a group, decide:
        is this possible? Defend your answer with the physics of circular
        motion --- where must the centre of any Earth orbit lie?
\end{enumerate}
\textit{(You should find $r \approx 4.2\times10^4$\,km, an altitude of about
36\,000\,km --- a tenth of the way to the Moon. And no: the centre of every
orbit must be the Earth's centre, so a satellite hovering off the equator is
impossible; it would need a circle whose centre is inside northern India.)}
\end{whiteboard}

\begin{supplementary}{The Tides, or: the Sea Breathes Twice a Day}
Bhutan is landlocked, so here is something many of us have never watched: at
the ocean's edge, the water level rises and falls by a metre or more, twice
every day. Newton's law explains this breathing --- and the explanation is
subtler than ``the Moon pulls the water.''

If the Moon simply pulled the water, there would be \emph{one} bulge of ocean
on the Moon's side, and one high tide per day as the Earth turns through it.
The observed rhythm --- two high tides a day --- says otherwise. The key is
that the Moon's pull is \emph{inverse-square}, so it is not uniform across
the Earth: the near-side ocean (about $59\,R_E$ from the Moon) is pulled
harder than the Earth's centre ($60\,R_E$), which in turn is pulled harder
than the far-side ocean ($61\,R_E$). Relative to the Earth, the near water
is pulled ahead, and the Earth is pulled away from the far water. Result:
\emph{two} bulges, one facing the Moon and one facing away, and every coast
passes through both each day. (The Sun pulls far harder in absolute terms,
but tides care about the \emph{difference} in pull across the Earth, and by
that measure the nearby Moon wins --- the Sun's tide is about half the
Moon's. When Sun and Moon line up, at new and full moon, the tides run
extra high.)

How big is the effect? In mid-ocean, under a metre. But geography can
amplify it spectacularly. The world record is the Bay of Fundy in eastern
Canada, where the water rises and falls up to \emph{sixteen metres} --- a
four-storey building, twice daily. The bay is a long, tapering funnel whose
natural sloshing period happens to be close to the 12.4-hour tidal rhythm,
so each tide arrives in step with the bay's own oscillation and pumps it
higher, the way well-timed pushes build up a playground swing. Resonance ---
an idea that will return, from swings to musical instruments to circuits.

Newton's tidal theory, published in the \emph{Principia} in 1687, was the
first time in history the tides had an explanation rather than a mere
timetable. Same law, third appearance: the apple, the Moon, and the sea.
\end{supplementary}
